\documentclass[aspectratio=169]{beamer}

\usetheme{Madrid}

\title{Educational Paxos Simulator}
\subtitle{Refactoring and Failure Scenarios}
\author{}
\date{}

\begin{document}

\begin{frame}
  \titlepage
\end{frame}

\begin{frame}{Presentation Summary}
  \begin{itemize}
    \item I extended my Paxos simulator after implementing the basic protocol.
    \item First, I refactored \texttt{simulator.py} to separate responsibilities.
    \item Then, I added \texttt{network.py} to simulate message delivery, delay, and loss.
    \item Finally, I implemented failure scenarios such as late old messages and leader conflict.
  \end{itemize}

  \begin{block}{Main goal}
    Understand why Paxos needs proposal numbers, Prepare / Promise, and careful message handling.
  \end{block}
\end{frame}

\begin{frame}{Refactored Simulator}
  \begin{block}{Before refactoring}
    \texttt{PaxosSimulator} handled too much logic by itself.
  \end{block}

  \begin{block}{After refactoring}
    The responsibilities are separated into \texttt{Proposer}, \texttt{Acceptor}, and \texttt{PaxosSimulator}.
  \end{block}

  \begin{itemize}
    \item \texttt{Proposer} creates messages and stores responses.
    \item \texttt{Acceptor} decides whether to promise or accept.
    \item \texttt{PaxosSimulator} only coordinates the protocol flow.
  \end{itemize}
\end{frame}

\begin{frame}{Why Refactoring Helps}
  \begin{itemize}
    \item The code now matches the Paxos roles more clearly.
    \item It is easier to understand what each component is responsible for.
    \item It prepares the project for network simulation.
    \item Failure scenarios can be implemented without putting all logic into the simulator.
  \end{itemize}

  \begin{block}{Key idea}
    The simulator should control the flow, but the Paxos logic should belong to the nodes.
  \end{block}
\end{frame}

\begin{frame}{Network.py}
  \begin{itemize}
    \item \texttt{network.py} simulates an unreliable network.
    \item It does not use real sockets or TCP.
    \item Instead, it manages messages inside the program.
  \end{itemize}

  \begin{block}{Main operations}
    \texttt{send()}, \texttt{delay()}, \texttt{drop()}, and \texttt{deliver\_one()}
  \end{block}

  \begin{itemize}
    \item This makes message delay and message loss explicit.
    \item It also makes Paxos failure scenarios easier to reproduce.
  \end{itemize}
\end{frame}

\begin{frame}{Late Old Message Scenario}
  \begin{enumerate}
    \item P1 sends \texttt{Prepare(1)}.
    \item P1 creates \texttt{AcceptRequest(1, A)}, but the network delays it.
    \item Before it arrives, P2 sends \texttt{Prepare(2)}.
    \item The Acceptor promises proposal number 2.
    \item Then the delayed \texttt{AcceptRequest(1, A)} arrives.
    \item The Acceptor rejects it because proposal number 1 is now old.
  \end{enumerate}
\end{frame}

\begin{frame}{Lesson from Late Old Message}
  \begin{block}{What this scenario shows}
    Old messages may arrive after newer messages in a distributed system.
  \end{block}

  \begin{itemize}
    \item Proposal numbers protect the system from old leaders.
    \item Once an Acceptor promises a newer proposal number, it rejects older requests.
    \item This prevents delayed old messages from breaking the current round.
  \end{itemize}
\end{frame}

\begin{frame}{Leader Conflict Scenario}
  \begin{itemize}
    \item This scenario intentionally skips \texttt{Prepare / Promise}.
    \item P1 sends \texttt{AcceptRequest(1, A)} to A1 and A2.
    \item A is chosen by a majority.
    \item P2 sends \texttt{AcceptRequest(2, B)} to A2 and A3.
    \item B is also chosen by a majority.
  \end{itemize}

  \begin{alertblock}{Problem}
    Two different values are chosen for the same slot.
  \end{alertblock}
\end{frame}

\begin{frame}{Why Paxos Prevents Conflict}
  \begin{itemize}
    \item In real Paxos, P2 cannot skip \texttt{Prepare}.
    \item P2 must collect \texttt{Promise} messages first.
    \item If a value was already accepted, the Promise includes that information.
    \item Therefore, P2 learns about the previous value before sending \texttt{AcceptRequest}.
  \end{itemize}

  \begin{block}{Key point}
    Prepare / Promise prevents a new leader from blindly overwriting a value that may already be chosen.
  \end{block}
\end{frame}

\begin{frame}{What I Learned}
  \begin{itemize}
    \item Paxos safety depends on both proposal numbers and majority agreement.
    \item Refactoring made the protocol roles easier to understand.
    \item Network simulation made message delay and message loss concrete.
    \item Leader conflict shows why Prepare / Promise is necessary.
    \item Paxos is easier to understand by implementing failure scenarios step by step.
  \end{itemize}
\end{frame}

\begin{frame}{Next Steps}
  \begin{itemize}
    \item Add more network-based scenarios.
    \item Implement network partition.
    \item Add safety tests.
    \item Extend the simulator from single-value Paxos to multi-slot Paxos.
    \item Write a README explaining each scenario.
  \end{itemize}
\end{frame}

\end{document}